% !TEX root =./ECE448_Practice_main.tex
%
% L7 practice set (AM and FM Receivers). Short by design: the L7 Lab packet
% carries the hands-on work. Blank copy: \begin{solution}[h] reserves h of work
% space. SOLUTIONS copy (\iskey=1) prints the red worked solutions.

\fullwidth{\textbf{Learning Objective 3.3: } I can implement an FM receiver in GNU
Radio, demodulate a live broadcast FM signal, and decode its RDS data.
\\[2pt]\textbf{Learning Objective 3.4: } I can explain and demonstrate the
trade-off between bandwidth and noise performance for AM vs.\ FM.
\\[4pt]\normalfont\small Show your work. Write a \textbf{Documentation} line at the
end (\textbf{None} if you did not collaborate).}

\begin{questions}

%%%%%%%%%%%%%%%%%%%%%%%%%%%%%%%%%%%%%%%%%
\question \textbf{FM bandwidth.} Carson's rule, a rule of thumb for setting a channel
filter, puts the width of an FM signal with peak deviation $\Delta f$ and audio
bandwidth $W$ at about
\eq{ B \approx 2(\Delta f + W) }
\begin{parts}

	\part Fill in the table. For the AM row, the station sends the same audio as DSB-AM.

	\begin{center}
	\renewcommand{\arraystretch}{1.25}
	\begin{tabular}{| l | c | c | p{1.4in} |}
	\hline
	Service & $\Delta f$ (kHz) & $W$ (kHz) & Bandwidth (kHz) \\\hline
	Broadcast FM, mono & 75 & 15 & \\\hline
	Two-way radio, NBFM & 5 & 3 & \\\hline
	AM, broadcast audio & -- & 15 & \\\hline
	\end{tabular}
	\end{center}

		\begin{solution}[0.9in]
		Broadcast FM: $2(75 + 15) = 180$ kHz. NBFM: $2(5 + 3) = 16$ kHz. AM: $2W = 30$ kHz.
		Broadcast FM is six times as wide as AM carrying the same audio.
		\end{solution}

	\part Open the FM Explorer on the Lesson 7 page and choose the broadcast preset,
	75~kHz of deviation with a 15~kHz tone. Lower the deviation to 25~kHz and describe
	what happens to the spectrum. Then explain in words why the broadcast signal is
	wider than $2\Delta f$ = 150~kHz.
		\begin{solution}[1.0in]
		At 25 kHz the lines crowd in toward the carrier and the shaded band shrinks to
		80 kHz. At 75 kHz the explorer draws lines every 15 kHz out past $\pm 75$ kHz:
		the frequency swings at the tone's rate, and that swing adds lines beyond the
		peak deviation, about one tone frequency on each side.
		\end{solution}

\end{parts}

\newpage
%%%%%%%%%%%%%%%%%%%%%%%%%%%%%%%%%%%%%%%%%
\question \textbf{The WBFM receiver.} A PlutoSDR records 2.4~MS/s centered on
98.8~MHz, and the station we want is at 99.1~MHz. The receiver is
\texttt{L07-wbfm-rx.grc}.
\begin{parts}

	\begin{minipage}[t]{0.58\linewidth}
		\part Find the center frequency of the Frequency Xlating FIR Filter, and the
		sample rate after it decimates by 10.
		\begin{solution}[0.9in]
		\eq{ 99.1 - 98.8 = 0.3\ \text{MHz} = 300\ \text{kHz} }
		\eq{ \dfrac{2.4\ \text{MS/s}}{10} = 240\ \text{kS/s} }
		\end{solution}
	\end{minipage}\hfill%
	\begin{minipage}[t]{0.37\linewidth}
		\raggedleft \vspace{12pt}
		shift = \ansbox[1.0in]{300 kHz}

		\vspace{4pt}
		rate = \ansbox[1.0in]{240 kS/s}
	\end{minipage}

	\part A classmate decimates by 20 instead. Explain what goes wrong, using the
	channel filter's ±90~kHz passband and the 57~kHz RDS subcarrier.
		\begin{solution}[1.0in]
		At 120 kS/s the flowgraph can hold only $-60$ to $+60$ kHz. The station's
		±90 kHz no longer fits, so its outer parts alias into the band, and after
		demodulation the RDS signal, which reaches 59.4 kHz, sits at the edge of the
		Nyquist frequency with no room for a filter.
		\end{solution}

	\begin{minipage}[t]{0.58\linewidth}
		\part Find the Quadrature Demod gain that makes 75~kHz of deviation produce an
		output of 1, and the output when the carrier sits 30~kHz above center.
		\begin{solution}[1.0in]
		\eq{ \dfrac{240\ \text{kHz}}{2\pi \times 75\ \text{kHz}} = 0.509 }
		\eq{ \dfrac{30}{75} = 0.40 }
		\end{solution}
	\end{minipage}\hfill%
	\begin{minipage}[t]{0.37\linewidth}
		\raggedleft \vspace{12pt}
		gain = \ansbox[1.0in]{$0.509$}

		\vspace{4pt}
		output = \ansbox[1.0in]{$0.40$}
	\end{minipage}

	\begin{minipage}[t]{0.58\linewidth}
		\part Find the decimation from the demodulator to a 48~kHz sound card, and the
		time constant to type into FM Deemphasis for a US station.
		\begin{solution}[1.0in]
		\eq{ \dfrac{240}{48} = 5 }
		75 $\mu$s in the US; Europe uses 50 $\mu$s.
		\end{solution}
	\end{minipage}\hfill%
	\begin{minipage}[t]{0.37\linewidth}
		\raggedleft \vspace{12pt}
		decimation = \ansbox[1.0in]{5}

		\vspace{4pt}
		$\tau$ = \ansbox[1.0in]{75 $\mu$s}
	\end{minipage}

	\part Explain why the flowgraph needs no Throttle, whether it plays a file or
	receives from the PlutoSDR.
		\begin{solution}[0.8in]
		The Audio Sink takes exactly 48,000 samples per second, so it paces a File
		Source. With the PlutoSDR, the radio's clock and the sound card's already set the
		pace, and a Throttle would add yet another clock and cause dropouts.
		\end{solution}

\end{parts}

\newpage
%%%%%%%%%%%%%%%%%%%%%%%%%%%%%%%%%%%%%%%%%
\question \textbf{RDS.} RDS sends 1187.5 bits per second on a 57~kHz subcarrier.
\begin{parts}

	\begin{minipage}[t]{0.58\linewidth}
		\part The RDS branch resamples to 19~kHz before Symbol Sync. Find the samples
		per bit that Symbol Sync is set for.
		\begin{solution}[0.8in]
		\eq{ \dfrac{19000}{1187.5} = 16 }
		\end{solution}
	\end{minipage}\hfill%
	\begin{minipage}[t]{0.37\linewidth}
		\raggedleft \vspace{12pt}
		\ansbox[1.0in]{16 samples}
	\end{minipage}

	\part The RDS signal is BPSK, a DSB-SC signal. Name the block in
	\texttt{L07-rds-rx.grc} that recovers its carrier, and explain why a recovered
	carrier $180^\circ$ off does no harm, given the differential coding.
		\begin{solution}[1.1in]
		The Constellation Receiver locks to the subcarrier's phase. A BPSK loop can lock
		at either of two phases $180^\circ$ apart, which inverts every bit. The data is
		differentially coded, so the Differential Decoder reads only whether each bit
		changed from the last, and inverting every bit leaves those changes the same.
		\end{solution}

	\part A station's audio is clear, but the station name never appears. Explain why
	RDS needs a stronger signal than the audio does.
		\begin{solution}[0.9in]
		The noise at the FM demodulator's output rises with frequency, so the 57 kHz
		subcarrier sees much more noise than the audio below 15 kHz, which de-emphasis
		also cleans up.
		\end{solution}

\end{parts}

\newpage
%%%%%%%%%%%%%%%%%%%%%%%%%%%%%%%%%%%%%%%%%
\question \textbf{AM and FM in the same noise.} \texttt{L07-am-vs-fm.grc} sends a
1~kHz tone as AM at $m = 1$ in a 10~kHz channel and as FM in a wider channel, with the
same transmitted power and the same noise. Its Input SNR slider sets the noise, and two
meters show each receiver's output SNR.
\begin{parts}

	\part Run it at 15~kHz deviation with Input SNR at 25~dB, as in the lesson figure.
	Say which output is cleaner, on the meters and on the frequency sink, and explain
	in words how FM wins although its wider channel lets in more noise.
		\begin{solution}[1.1in]
		FM, clearly: its meter reads well above AM's and its noise floor sits well under
		AM's. The demodulator's output follows how far the frequency swings, and a
		15 kHz swing is large next to the small jitter that the noise adds to a strong
		carrier. FM trades bandwidth for a cleaner output.
		\end{solution}

	\part Keep 15~kHz deviation, put the speaker on FM, and drag Input SNR down
	slowly. Describe what you hear, and what the FM meter does compared with the AM
	meter.
		\begin{solution}[1.1in]
		The tone stays clean at first. Then clicks start and get more frequent, and the
		FM meter falls quickly until it reads about the same as AM's. That collapse is
		the point where FM breaks into clicks. AM only hisses more as the slider comes
		down.
		\end{solution}

	\part You raise the deviation from 20 to 40~kHz. Say, in words, what happens to
	the FM meter while the signal is strong, and to where the clicks begin as the
	noise rises. Explain the trade.
		\begin{solution}[1.1in]
		With a strong signal the FM meter climbs. The clicks start sooner as the noise
		goes up, at a higher Input SNR. The wider deviation needs a wider channel, which
		lets in more noise, so the cleaner output lasts only while the signal is strong.
		\end{solution}

	\part Set Input SNR to 8~dB and listen to the FM output at 10, 20, and 40~kHz
	deviation. Say which sounds cleanest and why, in words. Then say which you would
	choose for a weak link to a distant receiver.
		\begin{solution}[1.1in]
		10 kHz sounds cleanest: a steady tone under hiss, a little better than AM, while
		20 kHz crackles with clicks and 40 kHz is mostly clicks. The narrow channel lets
		in the least noise, so at 10 kHz FM has not yet broken up. The wider deviations
		would sound cleaner on a strong link, but on a weak one the narrow deviation is
		the better choice.
		\end{solution}

\end{parts}

\newpage
%%%%%%%%%%%%%%%%%%%%%%%%%%%%%%%%%%%%%%%%%
\question \textbf{An AM receiver for ATIS.} An airport broadcasts its ATIS on
119.25~MHz. The PlutoSDR runs at 2.4~MS/s centered on 119.0~MHz, and the receiver is
\texttt{L07-am-rx.grc}.
\begin{parts}

	\begin{minipage}[t]{0.58\linewidth}
		\part Find the Frequency Xlating FIR Filter's center frequency, and the
		decimation that takes it straight to a 48~kHz sound card.
		\begin{solution}[0.9in]
		\eq{ 119.25 - 119.0 = 0.25\ \text{MHz} = 250\ \text{kHz} }
		\eq{ \dfrac{2.4\ \text{MS/s}}{48\ \text{kS/s}} = 50 }
		\end{solution}
	\end{minipage}\hfill%
	\begin{minipage}[t]{0.37\linewidth}
		\raggedleft \vspace{12pt}
		shift = \ansbox[1.0in]{250 kHz}

		\vspace{4pt}
		decimation = \ansbox[1.0in]{50}
	\end{minipage}

	\part Lesson 6 detected the envelope with Abs, a low-pass filter, and a gain of
	$\pi/2$. Explain why Complex to Mag replaces Abs and the $\pi/2$ gain.
		\begin{solution}[1.0in]
		After the Frequency Xlating FIR Filter, the station's carrier sits at 0 Hz, and
		each complex sample is the envelope times a phase factor. Its magnitude is the
		envelope itself, with no carrier to rectify or average away, so the receiver
		needs no $\pi/2$ correction. The low-pass filter stays, but it now serves only
		as the audio filter.
		\end{solution}

	\part Explain why the AM receiver needs an AGC and the FM receiver does not.
		\begin{solution}[1.0in]
		An envelope detector's output is proportional to the received amplitude, which
		changes with an aircraft's distance, so the AGC holds the audio level steady. The
		Quadrature Demod measures only how fast the phase turns, so the FM receiver's
		output does not depend on the signal's strength.
		\end{solution}

\end{parts}

\vspace{1.5em}
\noindent\textbf{Documentation:} \rule[-2pt]{0.6\linewidth}{0.4pt}

\end{questions}
